Vývoj \ac{APZ} výdajů v~ČR od vstupu do \ac{EU} (2004) v~porovnání s~nordickými zeměmi a~DE odhaluje, že integrace do \ac{EU} sama o~sobě nepřinesla konvergenci v~sociálních investicích: navzdory přístupu ke strukturálním fondům ČR nezvýšila \ac{APZ} výdaje v~podílu k~\ac{HDP} systémovým způsobem.
Epizodické výkyvy (vyšší výdaje v~krizových letech 2009--2010) svědčí o~reaktivním přístupu: \ac{APZ} je vnímána jako protikrizový nástroj, nikoliv jako permanentní investice do lidského kapitálu a~kvalitní pracovní trh.
Naproti tomu Dánsko a~Švédsko udržují vysoké \ac{APZ} výdaje i~v~době nízkézaměstnanosti: jde o~vyjádření filozofie, podle níž prevence a~kvalita trhu práce jsou prioritou bez ohledu na konjunkturální cyklus.
Zapojení odborů do správy \ac{APZ} programů --- jak je tomu v~nordickém systému --- je prvkem, který zvyšuje efektivitu výdajů: odbory mají zájem na úspěchu rekvalifikace (pracovníci si udrží zaměstnatelnost a~schopnost platit příspěvky), a~proto sociální dialog v~\ac{APZ} agendě je jak nástrojem efektivity, tak legitimity.

Dlouhodobý pohled potvrzuje strukturální podfinancování: ČR se dlouhodobě pohybuje v~pásmu \SIrange{0.2}{0.4}{\percent} \ac{HDP}, zatímco Dánsko investuje \SIrange{2.0}{2.5}{\percent} \ac{HDP} a~Německo \SIrange{0.7}{1.0}{\percent} \ac{HDP}.
Flexicurity model předpokládá trojici: flexibilní pracovní právo + štědré dávky v~nezaměstnanosti + masivní \ac{APZ}; ČR má v~zásadě jen první prvek, ostatní dva jsou výrazně podfinancovány.
Bez výrazného navýšení \ac{APZ} výdajů na alespoň \SI{1}{\percent} \ac{HDP} nemůže být flexibilizace trhu práce v~ČR legitimní: pracovníci nesou riziko flexibility, ale stát neposkytuje dostatečnou kompenzaci v~podobě přechodu a~rekvalifikace.

Vývoj \ac{APZ} výdajů v~ČR od vstupu do \ac{EU} (2004) v~porovnání s~nordickými zeměmi a~DE odhaluje, že integrace do \ac{EU} sama o~sobě nepřinesla konvergenci v~sociálních investicích: navzdory přístupu ke strukturálním fondům ČR nezvýšila \ac{APZ} výdaje v~podílu k~\ac{HDP} systémovým způsobem.
Epizodické výkyvy (vyšší výdaje v~krizových letech 2009--2010) svědčí o~reaktivním přístupu: \ac{APZ} je vnímána jako protikrizový nástroj, nikoliv jako permanentní investice do lidského kapitálu a~kvalitní pracovní trh.
Naproti tomu Dánsko a~Švédsko udržují vysoké \ac{APZ} výdaje i~v~době nízkézaměstnanosti: jde o~vyjádření filozofie, podle níž prevence a~kvalita trhu práce jsou prioritou bez ohledu na konjunkturální cyklus.
Zapojení odborů do správy \ac{APZ} programů --- jak je tomu v~nordickém systému --- je prvkem, který zvyšuje efektivitu výdajů: odbory mají zájem na úspěchu rekvalifikace (pracovníci si udrží zaměstnatelnost a~schopnost platit příspěvky), a~proto sociální dialog v~\ac{APZ} agendě je jak nástrojem efektivity, tak legitimity.

Dlouhodobý pohled potvrzuje strukturální podfinancování: ČR se dlouhodobě pohybuje v~pásmu \SIrange{0.2}{0.4}{\percent} \ac{HDP}, zatímco Dánsko investuje \SIrange{2.0}{2.5}{\percent} \ac{HDP} a~Německo \SIrange{0.7}{1.0}{\percent} \ac{HDP}.
Flexicurity model předpokládá trojici: flexibilní pracovní právo + štědré dávky v~nezaměstnanosti + masivní \ac{APZ}; ČR má v~zásadě jen první prvek, ostatní dva jsou výrazně podfinancovány.
Bez výrazného navýšení \ac{APZ} výdajů na alespoň \SI{1}{\percent} \ac{HDP} nemůže být flexibilizace trhu práce v~ČR legitimní: pracovníci nesou riziko flexibility, ale stát neposkytuje dostatečnou kompenzaci v~podobě přechodu a~rekvalifikace.

Vývoj \ac{APZ} výdajů v~ČR od vstupu do \ac{EU} (2004) v~porovnání s~nordickými zeměmi a~DE odhaluje, že integrace do \ac{EU} sama o~sobě nepřinesla konvergenci v~sociálních investicích: navzdory přístupu ke strukturálním fondům ČR nezvýšila \ac{APZ} výdaje v~podílu k~\ac{HDP} systémovým způsobem.
Epizodické výkyvy (vyšší výdaje v~krizových letech 2009--2010) svědčí o~reaktivním přístupu: \ac{APZ} je vnímána jako protikrizový nástroj, nikoliv jako permanentní investice do lidského kapitálu a~kvalitní pracovní trh.
Naproti tomu Dánsko a~Švédsko udržují vysoké \ac{APZ} výdaje i~v~době nízkézaměstnanosti: jde o~vyjádření filozofie, podle níž prevence a~kvalita trhu práce jsou prioritou bez ohledu na konjunkturální cyklus.
Zapojení odborů do správy \ac{APZ} programů --- jak je tomu v~nordickém systému --- je prvkem, který zvyšuje efektivitu výdajů: odbory mají zájem na úspěchu rekvalifikace (pracovníci si udrží zaměstnatelnost a~schopnost platit příspěvky), a~proto sociální dialog v~\ac{APZ} agendě je jak nástrojem efektivity, tak legitimity.

Dlouhodobý pohled potvrzuje strukturální podfinancování: ČR se dlouhodobě pohybuje v~pásmu \SIrange{0.2}{0.4}{\percent} \ac{HDP}, zatímco Dánsko investuje \SIrange{2.0}{2.5}{\percent} \ac{HDP} a~Německo \SIrange{0.7}{1.0}{\percent} \ac{HDP}.
Flexicurity model předpokládá trojici: flexibilní pracovní právo + štědré dávky v~nezaměstnanosti + masivní \ac{APZ}; ČR má v~zásadě jen první prvek, ostatní dva jsou výrazně podfinancovány.
Bez výrazného navýšení \ac{APZ} výdajů na alespoň \SI{1}{\percent} \ac{HDP} nemůže být flexibilizace trhu práce v~ČR legitimní: pracovníci nesou riziko flexibility, ale stát neposkytuje dostatečnou kompenzaci v~podobě přechodu a~rekvalifikace.

Vývoj \ac{APZ} výdajů v~ČR od vstupu do \ac{EU} (2004) v~porovnání s~nordickými zeměmi a~DE odhaluje, že integrace do \ac{EU} sama o~sobě nepřinesla konvergenci v~sociálních investicích: navzdory přístupu ke strukturálním fondům ČR nezvýšila \ac{APZ} výdaje v~podílu k~\ac{HDP} systémovým způsobem.
Epizodické výkyvy (vyšší výdaje v~krizových letech 2009--2010) svědčí o~reaktivním přístupu: \ac{APZ} je vnímána jako protikrizový nástroj, nikoliv jako permanentní investice do lidského kapitálu a~kvalitní pracovní trh.
Naproti tomu Dánsko a~Švédsko udržují vysoké \ac{APZ} výdaje i~v~době nízkézaměstnanosti: jde o~vyjádření filozofie, podle níž prevence a~kvalita trhu práce jsou prioritou bez ohledu na konjunkturální cyklus.
Zapojení odborů do správy \ac{APZ} programů --- jak je tomu v~nordickém systému --- je prvkem, který zvyšuje efektivitu výdajů: odbory mají zájem na úspěchu rekvalifikace (pracovníci si udrží zaměstnatelnost a~schopnost platit příspěvky), a~proto sociální dialog v~\ac{APZ} agendě je jak nástrojem efektivity, tak legitimity.

Dlouhodobý pohled potvrzuje strukturální podfinancování: ČR se dlouhodobě pohybuje v~pásmu \SIrange{0.2}{0.4}{\percent} \ac{HDP}, zatímco Dánsko investuje \SIrange{2.0}{2.5}{\percent} \ac{HDP} a~Německo \SIrange{0.7}{1.0}{\percent} \ac{HDP}.
Flexicurity model předpokládá trojici: flexibilní pracovní právo + štědré dávky v~nezaměstnanosti + masivní \ac{APZ}; ČR má v~zásadě jen první prvek, ostatní dva jsou výrazně podfinancovány.
Bez výrazného navýšení \ac{APZ} výdajů na alespoň \SI{1}{\percent} \ac{HDP} nemůže být flexibilizace trhu práce v~ČR legitimní: pracovníci nesou riziko flexibility, ale stát neposkytuje dostatečnou kompenzaci v~podobě přechodu a~rekvalifikace.

\input{texparts/python/lmp_expenditure_2004}



